\section{Background and Related Work}
\label{sec:related}

\subsection{The Vedic corpus as a computational object}

The four Vedic Sa\d{m}hit\=as are among the oldest continuously transmitted texts in
any language, preserved orally with a fidelity that makes them unusually good
material for textual computation and unusually unforgiving of sloppy addressing. Each
survives in named \iast{\'s\=akh\=a}s whose verse inventories differ; a claim about
``the Atharvaveda'' that does not say \'Saunaka or Paippal\=ada is not yet a claim
about a corpus. The standard modern scholarly translation of the \d{R}gveda
\citep{jamison2014rigveda} is explicit about this, and the nineteenth-century
translations that remain the only complete English renderings reachable for parts of
the corpus \citep{griffith1896rigveda,whitney1905atharva} each carry their own
numbering conventions, which is a recurring source of alignment error
(\cref{sec:translation}).

One reference work deserves particular mention because it is the direct predecessor
of a layer reported here. Bloomfield's \emph{Vedic Concordance}
\citep{bloomfield1906concordance} is an alphabetic index to every line of every
stanza of the published Vedic literature, with an account of the variations across
Vedic books --- that is, a manual, exhaustive index of exactly the cross-textual
repetition that \cref{sec:connections} detects computationally. We did not use it as
a gold standard, because the project does not hold a machine-readable form of it
under rights permitting that use; we note the omission as the most obvious missing
evaluation in this paper (\cref{sec:limitations}).

For the S\=amaveda specifically, the relationship between the \=Arcika verse text
and its sung realisation is the subject of a substantial musicological literature
\citep{howard1977samavedic}. That literature is the reason \cref{sec:case-samaveda}
is careful: the project records svara marks as printed and declines to decode them
into pitch, and the decipherment authority it would need in order to do so
responsibly is not held.

\subsection{Sanskrit computational linguistics}

The computational infrastructure for Sanskrit has matured substantially. The Digital
Corpus of Sanskrit \citep{hellwig2021dcs} supplies morphological and lexical
annotation at scale; neural segmentation now handles the joint problem of compound
splitting and sandhi resolution without hand-built features
\citep{hellwig2018segmentation}; the Treebank of Vedic Sanskrit
\citep{hellwig2020vedic} supplies human-validated dependency annotation over Vedic
material specifically; and VedaWeb \citep{kiss2019vedaweb} provides a research
platform over morphologically and metrically annotated Old Indo-Aryan text. A
digital Anukrama\d{n}\=\i\ \citep{akavarapu2023anukramani} makes the traditional
hymn-level index --- seer, deity, metre --- machine-readable for the \d{R}gveda.

VedGraph consumes rather than extends these. The treebank is the external annotation
against which its reference set is constructed (\cref{sec:gold}); the Anukrama\d{n}\=\i\
is the source of its \textsc{l1} attribution layer; VedaWeb supplies a parallel
reading and a linguistic annotation layer. The contribution claimed here is not a new
language resource but an architecture for holding claims of very different
epistemic kinds in one queryable structure.

\subsection{Knowledge-graph construction and refinement}

Surveys of knowledge graphs \citep{hogan2021knowledge} and of graph refinement
\citep{paulheim2017refinement} establish the vocabulary this paper uses and also its
central gap: refinement is generally framed as improving \emph{correctness} and
\emph{completeness}, with provenance treated as metadata about a dataset rather than
as a filterable property of each assertion. Where per-claim provenance is modelled,
the established mechanisms are reification and its successors --- nanopublications
\citep{groth2010nanopublication}, which package an assertion with its provenance and
publication info as named graphs, and PROV-O \citep{lebo2013provo}, which supplies a
general vocabulary for entities, activities and agents.

VedGraph's model is a considerably less expressive special case of what those
standards can express, and deliberately so. Rather than a general provenance graph
it stores four fixed, closed-vocabulary properties on every edge
(\cref{sec:model}), which buys two things a general model does not: the axes are
indexed, so ``show me only source-explicit claims'' is cheap rather than a join; and
a closed vocabulary admits universal-quantifier invariants, so a missing or
out-of-enum value is a build failure rather than an under-specified record. The
trade is expressiveness for enforceability. We would expect a PROV-O export to be
straightforward and have not built one.

\subsection{Language models in extraction, and their failure modes}

The programme of using language models to build and complete knowledge graphs is
well surveyed \citep{pan2024unifying}. The failure mode this paper is organised
around is equally well documented: generated text can be fluent, confident and
unsupported \citep{ji2023hallucination}, and --- crucially for any design that hopes
to let a model check itself --- models do not reliably self-correct reasoning
without external feedback, with performance sometimes degrading after a
self-correction pass \citep{huang2024selfcorrect}.

That last result concerns \emph{intrinsic} self-correction --- a model revising its
own answer with no external feedback --- which is not the regime described here: the
attacking agent carries a different instruction and has read access to the store,
which is the external-feedback condition that work excludes. It is therefore a
motivation for the division of labour reported here rather than evidence that the
division works, and we are aware of no result that predicts it will.
VedGraph does not ask a model to assess its own output; it asks a
\emph{differently-instructed} agent to attack that output with a stated hypothesis,
and it requires the attack to be expressed as a measurement against the store rather
than as a judgement (\cref{sec:multiagent}). We are explicit in
\cref{sec:threats} that this does not achieve statistical independence, and that a
correlated failure across a shared model family would be invisible to the whole
regime.

On agent architectures, interleaving reasoning with tool-use actions
\citep{yao2023react} is the pattern the construction runtime instantiates, but
\emph{only} as a pattern: no orchestration framework was used, and we verified that
by inspecting the dependency manifests (\cref{sec:appendix-agentic}). We mention the
literature to place the work, not to claim a framework the repository does not
contain.

\subsection{Text reuse detection}

Detecting reuse in historical corpora is an established line with mature tools:
scalable passage-cluster detection over large noisy collections
\citep{smith2013infectious}, intertext detection tuned for classical poetry
\citep{coffee2013tesserae}, and a general algorithm suite for historical reuse
\citep{buechler2014tracer}. Our candidate-generation strategy --- banded MinHash over
character $n$-grams, then a composite score --- is conventional against that
background, and \cref{sec:near} reports the parameters rather than claiming novelty
for the method.

Two things here are less conventional. First, the comparison \emph{surface} is
treated as a research object rather than as preprocessing: six named match levels
are kept distinct and the code records which one a given edge was matched on,
because a script-folded match and a byte-identical match assert different things
(\cref{sec:ladder}). Second, \emph{direction}. The reuse literature generally
establishes direction from external chronology. Because Vedic relative chronology is
contested, this project measures a structural signal instead --- compilation
asymmetry (\cref{sec:reuse}) --- and validates it per scope, accepting it where it
recovers a tradition's own claim at 0.991 consistency and refusing it where it
cannot separate the sides.

The corpus's formulaic character connects to oral-formulaic theory
\citep{lord1960singer}, which supplies the reason a formula layer is worth building
at all: in an orally composed and orally transmitted corpus, a repeated multi-word
span is a unit of composition rather than an accident. We use that as motivation and
make no claim to contribute to the theory.

\subsection{Evaluation without a gold standard}

Reporting agreement figures properly requires care about what the coefficient
measures and what population it was computed over \citep{artstein2008agreement}. The
situation in this project is more constrained than the usual one: there is no human
annotation of any VedGraph claim at all, so no inter-annotator coefficient over its
assertions is computable even in principle (\cref{sec:gold}). What exists is an
external, human-validated annotation of the \emph{text} \citep{hellwig2020vedic},
against which claims can be adjudicated mechanically. The project names that
arrangement with its own term, and refuses to call it gold. Everything we report
about the semantic layer's accuracy is bounded by that.

Finally, the data-management principles the project's registries implement --- rights
and provenance recorded per artefact, identifiers stable and resolvable, machine-
readable scope --- are those of the FAIR framework \citep{wilkinson2016fair}, with
the important qualification that reusability here is constrained by third-party
rights the project does not hold (\cref{sec:data-availability}).
